 \pdfoutput=1
\documentclass[11pt]{article}

% -------------------------------------------------------
% Core packages
% -------------------------------------------------------
\usepackage[T1]{fontenc}
\usepackage[utf8]{inputenc}
\usepackage[margin=1.1in]{geometry}
\usepackage{amsmath, amssymb, amsthm}
\usepackage{mathtools}
\usepackage[final,expansion=false]{microtype}

% -------------------------------------------------------
% Figures and tables
% -------------------------------------------------------
\usepackage{graphicx}
\usepackage{caption}
\usepackage{subcaption}
\usepackage{float}
\usepackage{booktabs}
\usepackage{array}

% -------------------------------------------------------
% Colors and boxes
% -------------------------------------------------------
\usepackage{xcolor}
\usepackage{mdframed}

\definecolor{boxblue}{RGB}{235,243,252}
\definecolor{boxborder}{RGB}{100,160,220}
\definecolor{greenbg}{RGB}{235,250,238}
\definecolor{greenborder}{RGB}{60,170,90}
\definecolor{warnbg}{RGB}{255,248,230}
\definecolor{warnborder}{RGB}{220,160,30}
\definecolor{graybg}{RGB}{245,245,245}
\definecolor{grayborder}{RGB}{180,180,180}

\newmdenv[backgroundcolor=boxblue,linecolor=boxborder,linewidth=1.2pt,
  innerleftmargin=12pt,innerrightmargin=12pt,innertopmargin=10pt,
  innerbottommargin=10pt,roundcorner=4pt]{bluebox}
\newmdenv[backgroundcolor=greenbg,linecolor=greenborder,linewidth=1.2pt,
  innerleftmargin=12pt,innerrightmargin=12pt,innertopmargin=10pt,
  innerbottommargin=10pt,roundcorner=4pt]{greenbox}
\newmdenv[backgroundcolor=warnbg,linecolor=warnborder,linewidth=1.2pt,
  innerleftmargin=12pt,innerrightmargin=12pt,innertopmargin=10pt,
  innerbottommargin=10pt,roundcorner=4pt]{warnbox}
\newmdenv[backgroundcolor=graybg,linecolor=grayborder,linewidth=0.8pt,
  innerleftmargin=12pt,innerrightmargin=12pt,innertopmargin=10pt,
  innerbottommargin=10pt,roundcorner=4pt]{graybox}

% -------------------------------------------------------
% References — plain biblio, arXiv safe
% -------------------------------------------------------
\usepackage{cite}
\usepackage[colorlinks=true,
            linkcolor=blue!60!black,
            citecolor=blue!60!black,
            urlcolor=blue!60!black]{hyperref}
\usepackage{url}

% -------------------------------------------------------
% Misc
% -------------------------------------------------------
\usepackage{parskip}
\usepackage{enumitem}
\usepackage{algorithm}
\usepackage{algpseudocode}

% -------------------------------------------------------
% Theorem environments — AMS style
% -------------------------------------------------------
\theoremstyle{plain}
\newtheorem{theorem}{Theorem}[section]
\newtheorem{lemma}[theorem]{Lemma}
\newtheorem{proposition}[theorem]{Proposition}
\newtheorem{corollary}[theorem]{Corollary}
\newtheorem{assumption}[theorem]{Assumption}

\theoremstyle{definition}
\newtheorem{definition}{Definition}[section]
\newtheorem{example}{Example}[section]

\theoremstyle{remark}
\newtheorem*{remark}{Remark}
\newtheorem*{notation}{Notation}

% -------------------------------------------------------
% Custom commands
% -------------------------------------------------------
\newcommand{\Gbar}{\bar{G}}
\newcommand{\E}{\mathbb{E}}
\newcommand{\R}{\mathbb{R}}
\newcommand{\dW}{dW_t}
\newcommand{\pV}[1]{\frac{\partial V}{\partial #1}}
\newcommand{\ppV}[1]{\frac{\partial^2 V}{\partial #1^2}}
\newcommand{\norm}[1]{\left\|#1\right\|}
\newcommand{\abs}[1]{\left|#1\right|}

% -------------------------------------------------------
% Title block — arXiv format
% -------------------------------------------------------
\title{
  \textbf{Optimal Bus Holding via Stochastic Control:} \\
  \textbf{An Empirical Hamilton--Jacobi--Bellman Approach} \\[0.5em]
  \large Calibrated to MTA M15 GPS Data
}

\author{
  Franco \\[0.3em]
  \small Department of Statistics, University of Central Florida \\
  \small \texttt{your.email@ucf.edu}
}

\date{\today}

% -------------------------------------------------------
\begin{document}
\maketitle

% -------------------------------------------------------
% MSC 2020 classifications — required for math.OC
% -------------------------------------------------------
\vspace{-0.5em}
\noindent\textbf{MSC 2020:}
49L20 (Dynamic programming in optimal control),
60H10 (Stochastic ordinary differential equations),
93E20 (Optimal stochastic control),
90B20 (Traffic problems)

\medskip
\noindent\textbf{Keywords:}
stochastic optimal control,
Hamilton--Jacobi--Bellman equation,
bus bunching,
Ornstein--Uhlenbeck process,
dynamic programming,
finite difference methods

\bigskip

% -------------------------------------------------------
\begin{abstract}
% 150-250 words. Cover: problem, model, method, data, result.
% Structure: (1) motivation, (2) model, (3) method, (4) data, (5) result.
\end{abstract}

\tableofcontents
\newpage

% -------------------------------------------------------
\section{Introduction}
\label{sec:intro}

% Paragraphs:
% 1. Bus bunching problem and why it matters
% 2. Existing approaches and their limitations
% 3. Our approach: stochastic control + real data
% 4. Finance parallel (Merton, pairs trading)
% 5. Paper outline

% -------------------------------------------------------
\section{Problem Formulation}
\label{sec:formulation}

\subsection{State Variable and Dynamics}
\label{sec:dynamics}

\subsection{Control Variable}
\label{sec:control}

\subsection{Objective Function and Value Function}
\label{sec:objective}

% -------------------------------------------------------
\section{The Hamilton--Jacobi--Bellman Equation}
\label{sec:hjb}

\subsection{Bellman's Principle of Optimality}
\label{sec:bellman}

\subsection{It\^{o}'s Lemma and HJB Derivation}
\label{sec:ito}

\subsection{The HJB Equation}
\label{sec:hjb_eq}

\subsection{Optimal Bang-Bang Policy}
\label{sec:bangbang}

% -------------------------------------------------------
\section{Numerical Solution}
\label{sec:numerics}

\subsection{Spatial and Temporal Discretization}
\label{sec:grid}

\subsection{Finite Difference Approximations}
\label{sec:fd}

\subsection{Explicit Backward Sweep}
\label{sec:sweep}

\subsection{Stability: CFL Conditions}
\label{sec:cfl}

% -------------------------------------------------------
\section{Empirical Calibration}
\label{sec:data}

\subsection{Data Source and Preprocessing}
\label{sec:datasource}

\subsection{OU Parameter Estimation via AR(1) Regression}
\label{sec:estimation}

\subsection{Fitted Parameters}
\label{sec:params}

% -------------------------------------------------------
\section{Results}
\label{sec:results}

\subsection{Value Function}
\label{sec:vf_results}

\subsection{Optimal Holding Policy Surface}
\label{sec:policy_results}

\subsection{Monte Carlo Evaluation}
\label{sec:mc_results}

% -------------------------------------------------------
\section{Connection to Quantitative Finance}
\label{sec:finance}

\subsection{Structural Equivalence to Merton's Problem}
\label{sec:merton}

\subsection{Pairs Trading Analogy}
\label{sec:pairs}

% -------------------------------------------------------
\section{Discussion}
\label{sec:discussion}

\subsection{Limitations}
\label{sec:limitations}

\subsection{Extensions}
\label{sec:extensions}

% -------------------------------------------------------
\section{Conclusion}
\label{sec:conclusion}

% -------------------------------------------------------
\appendix

\section{Derivation of the HJB Equation}
\label{app:hjb}

\section{AR(1) to OU Parameter Conversion}
\label{app:ou}

\section{Finite Difference Derivations from Taylor Series}
\label{app:fd}

\section{CFL Stability Analysis}
\label{app:cfl}

% -------------------------------------------------------
\bibliographystyle{plain}
\bibliography{references}

\end{document}
